\section{Discussion}
\label{sec:discussion}

The result (\S\ref{sec:results}) is neither cleanly ``portable'' nor
cleanly ``fragile'' -- it is source-pair-dependent in a specific,
mechanistically legible way, and that dependence itself is the finding.

\subsection{Portability Is Real but Source-Pair-Dependent, Not Uniform}

Comparative scheduler rankings are never uncorrelated between any two
sources at any region (\S\ref{sec:cross-source-rankings}): a scheduler
comparison performed on one of this study's three sources retains
positive aggregate ranking agreement with the others. At the same time,
portability is far from uniform, and the non-uniformity is legible rather
than diffuse: Azure-2024 and Bailian/Qwen agree almost perfectly at every
operating region, while every comparison involving BurstGPT is consistently
lower (\S\ref{sec:cross-source-rankings},
\S\ref{sec:results-robustness}). The
practical implication is neither ``a small, well-chosen evaluation
footprint always suffices'' nor ``every scheduler comparison must be
re-verified on every workload'': it is that \emph{which} second source a
comparison is checked against matters more than \emph{whether} a second
source is checked at all. A comparison validated only on Azure-2024-like
traffic shows strong agreement with Bailian/Qwen-like traffic under
this study's panel and metric, but the same transfer to BurstGPT-like
traffic is measurably less reliable, and that reliability gap holds at
every load level tested, not only under stress.

\subsection{Practically Supported Reversals Are Real, Uncommon, and
Load-Concentrated}

Practically and statistically supported reversals are a small minority of
PRIMARY-metric pairwise comparisons (\S\ref{sec:reversals}), but not zero,
and not dismissible as measurement noise: the same bootstrap procedure
classifies the large majority of comparisons as stable. Every supported
reversal we examined is the same policy pair and involves BurstGPT on one
side, and reversals are not spread evenly across the operating envelope;
they concentrate at and above \texttt{PRE\_KNEE}. A scheduler comparison performed only at low load,
even on a source pair that is otherwise highly portable, would have missed
every reversal this benchmark found. This motivates evaluating scheduler
comparisons across an operating curve, not at a single convenient load
point, as this benchmark's design already does
(\S\ref{sec:calibration}).

\subsection{Practical Implications for Benchmark Design}

Two implications follow directly and are not deployment recommendations
about which scheduler to use, only about how to evaluate schedulers
credibly. First, a comparison checked against only one additional workload
source cannot be assumed representative of ``other workloads'' in
general -- which specific source is used as the check matters, and
BurstGPT-like traffic in particular is where this study's evidence
concentrates its cross-source disagreement. Second, benchmark sample
complexity is itself load- and metric-behavior-relevant evidence:
\S\ref{sec:sample-complexity}'s finding that exact-order recovery requires
the full window budget for two of three sources, while top-1 recovery is
achievable from very few windows for those same two sources, means a
researcher who only needs to know ``which scheduler is best'' faces a
different, much smaller evidence requirement than one who needs the
complete ordering -- a distinction a single point-estimate benchmark
result cannot convey on its own.

\subsection{Portability Also Depends on Which Metric and Which SLO
Definition Is Used}

\S\ref{sec:cross-source-rankings}--\S\ref{sec:reversals} establish that a
ranking's portability depends on which \emph{workload source} it is checked
against. \S\ref{sec:results-cross-metric} and \S\ref{sec:results-slo} show
that portability depends, independently, on which \emph{metric} and which
\emph{SLO definition} is used to compute the ranking in the first place --
two further conditioning variables, not a repetition of the source-pair
finding above.

The cross-metric result is the more consequential of the two. Rank
correspondence across metric pairs is only moderate (median
$\tau_b = 0.419$) -- well below the agreement between the two most portable
\emph{sources} -- and metrics sometimes rank policies in opposite order
rather than merely disagreeing in degree. Across the evaluated metric pairs
within source-region conditions, the top-ranked policy differs more often
than it agrees (\S\ref{sec:results-cross-metric}).
This does not mean
metrics are arbitrary or that any one choice is as good as another -- ANWG,
completion fraction, and the latency- and SLO-based metrics measure
genuinely different deployment objectives (\S\ref{sec:metrics}), and a
practitioner optimizing for tail latency has no reason to expect the same
ranking as one optimizing for weighted goodput. The finding is instead that
\emph{which} objective a benchmark reports on is as much a part of the
ranking-portability question as which workload it was measured on, and a
scheduler comparison that reports only its primary metric is silent on
whether that comparison would hold under a deployment's actual objective.

The SLO-definition result is distinct from metric sensitivity and
comparatively reassuring, but not a clean pass: ranking robustness under
alternative SLO-synthesis conventions is closer to the cross-source ceiling
than to the cross-metric result, yet the top-ranked policy and a non-trivial
minority of specific reversals do change with the SLO definition
(\S\ref{sec:results-slo}). The practical reading is narrower than either
``robust'' or ``fragile'': a scheduler conclusion can remain broadly
ordered under a nearby SLO definition even while its winning policy, or a
specific reversal a reader cares about, does not.

Taken together with \S\ref{sec:cross-source-rankings}'s source-pair finding
and \S\ref{sec:reversals}'s load-concentration finding, this study's central
conclusion is that a comparative scheduler claim is conditional on at least
four axes -- workload source, operating region, evaluation metric, and (for
deadline-dependent policies) SLO definition -- and reporting a ranking
without stating which values of each axis produced it risks the ranking
being read as more general than the evidence supports.

\subsection{Explanatory Telemetry: A Preliminary, Heavily Qualified Read}
\label{sec:telemetry-explanation}

The frozen, pre-specified explanatory model (\S\ref{sec:telemetry})
associates the five workload
descriptors with reversal sites at each (policy pair, condition) cell, but
converges with a usable sample size ($\geq 20$ observations) for only 42 of
990 cells -- most reversal sites are too sparse for this per-cell model to
fit reliably, and the frozen output records point coefficients only, with
no standard errors or significance test. Any reading of this evidence must
therefore be provisional. Within that small, non-significance-tested set,
\texttt{prompt\_tokens\_cv} has a negative coefficient in 39 of 42 cells --
the most consistent single direction among the five descriptors -- which
is at least \emph{consistent with} (not confirmatory of) prompt-length
variability being associated with where reversals occur, echoing
\S\ref{sec:drivers}'s finding that prompt-length structure is the
classifier's most predictive, non-redundant feature for source
separability itself. The other four descriptors show
weaker, closer-to-even sign splits across the interpretable cells and do
not support a directional claim. We do not read this as evidence that any
descriptor \emph{causes} a reversal: because the sources are themselves
observably separable (\S\ref{sec:separability}), an association between a
descriptor and a reversal may reflect only a source-separating feature that
co-occurs with where reversals were or were not observed. It is a candidate
association worth a properly powered, purpose-built follow-up study;
consistent with \S\ref{sec:evidence-boundaries}, this model is explanatory
only and is never used to select, route, or recommend a scheduler.

\subsection{Physical-Execution Evidence}

\S\ref{sec:real-system}'s selected-case physical validation is informative
primarily about a boundary of simulator fidelity, not as an independent
statistical estimate of how often this benchmark's simulated findings
would hold on real hardware in general (scope limits:
\S\ref{sec:threats}). The non-reproduced reversal exposes a fidelity
limitation specific to that selected comparison, while the stable control
remained qualitatively stable -- consistent with, but not a confirmation
of, the simulator's broader findings. Together they show the same
case-dependence documented in \S\ref{sec:results}, now at the
simulator-to-hardware boundary, without adding statistical weight to those
earlier findings.

\subsection{Benchmark Construction and Practical Reproducibility}

Regardless of outcome direction, this study's paired design, fixed
scheduler panel, and preregistered protocol (\S\ref{sec:study-design})
are intended to make the benchmark itself reusable: a future scheduler
design can be dropped into the same panel and evaluated against the same
frozen windows, operating regions, and metric contract, using the public
artifact (\S\ref{sec:artifact}) rather than re-deriving the protocol
from this manuscript's prose.
